\documentclass[title=Monsieur, name=Gaboué, r=0.62, g=0.13, b=0.44]{../mrms}
\begin{document}

\coverpage{../img/m-gaboue.jpg}
\collectionpage
\titlepage

% Monsieur Gaboué se réveille tout gargouillant et croustillant
% Ça fait des mois qu'il procrastine le trou dans son mur, mais là il fait froid (on est l'automne)
% Il veut faire une partie de Brun, mais il a trop froid et faim
% Il va pour se faire à manger, mais il n'y a que des lentilles, et il a envie d'haricots, alors il va faire une "petite" marche pour se chercher à manger au restaurant "Chez Matonrière"
% Sur le chemin, il croise monsieur médiocre, qui lui dit que grâce à son château, il est au moins à moitié son dieu, alors monsieur Gaboué lui fait un finger
% Sur le chemin, il croise mme fin de phrase, qui n'arrive pas à finir de dire quoi que ce soit, alors monsieur Gaboué lui fait un finger
% chez matonrière, il tombe sur son ami monsieur arrêter-quand-je-veux, qui lui dit qu'il connait quelqu'un qui pourrait l'aider avec son problème de mur
% repartant avec son bol de binnes, il tombe sur mme nata, mais c'est une tartelette alors elle dit rien, alors monsieur gaboué lui fait un finger
% il arrive chez monsieur ça-fait-la-job, qui peut lui faire ça _cheap_, puis qui lui refile un mur de tôle. Monsieur gaboué a pas envie de l'installer lui-même
% monsieur gaboué continue son chemin jusque chez lui, il arrive pis m ça-fait-la-job est déjà là (y'ont pris le bus)
% c'est le chaos
% dans le chaos, un mur de tôle se fait installer
% gaboué pogne les nerfs pis kickout tout le monde
% pis maintenant que tout est chill, il appelle son ami m arrêter-quand-je-veux... pour faire une partie de brun
% la musique micro-tonale c'est micro-louche

\txtpage{
    Monsieur Gaboué entrouvrit les yeux. Criiiiic.

    Ses yeux étaient tout croustillants: Il faisait très froid chez lui.

    Après un effort considérable, il ouvrit les yeux, et regarda son bon vieux trou-dans-le-mur.
}
\txtpage{<closeup de gaboué qui ouvre ses yeux?>}

\txtpage{
    Il y a quelques mois, le mur s'était subitement écroulé, alors que le sol de la rue s'était
    abaissé pour aucune bonne raison. Maintenant que l'automne était arrivé, monsieur
    Gaboué ne pouvait plus se permettre de l'ignorer.

    Avant, il devait faire quelques chose à propos de son ventre gargouillant.
}
\txtpage{<trou dans le mur, avec la rue toute croche?>}

\txtpage{
    Monsieur Gaboué jeta un coup d'oeil dans sa cuisine pour quelque chose à manger, mais tout ce qu'il y
    avait c'était des lentilles. Souvent, les lentilles c'est très acceptable, mais là monsieur Gaboué
    se dit que ça prendrait quelque chose de plus réconfortant, vu la température.

    Il partit donc de chez lui, pour faire une petite marche de quelques dizaines de kilomètres
    jusqu'au restaurant "Chez Matonrière".
}
\txtpage{<Pot de lentilles>}

\txtpage{
    Au bout de quelques minutes de marche, une voix ténébreuse retentit:

    - J'ai un château, alors je suis au moins à demi ton dieu! Sois exulté en ma présence!
    Exigea monsieur Médiocre.

    - Ça dépend du château. Tu as une photo? Demanda monsieur Gaboué.

    Monsieur Médiocre lui montra la photo.

    Monsieur Gaboué fronça les sourcils, rendit son jugement, puis répondit à monsieur Médiocre d'un simple geste avant de repartir.
}
\txtpage{<Finger à monsieur médiocre, monsieur médiocre tient une photo>}

\txtpage{
    Plus loin, monsieur Gaboué croisa madame Fin-De-Phrase.

    - Monsieur Gaboué, ça tombe à pic! Avec monsieur Arrêter-Quand-Je-Veux, nous organisons
    une... une...

    - Fin de phrase? compléta monsieur Gaboué.

    - Ha. Fit sèchement madame Fin-De-Phrase. Elle reprit: Nous organisons une partie de ... de ...

    Monsieur Gaboué avait compris, mais il avait aussi trop faim, alors il lui dit au-revoir d'un simple geste puis repris son chemin.
}
\txtpage{<Finger à mme fin-de-phrase>}

\txtpage{
    Quelques heures plus tard, un monsieur Gaboué affamé arriva au restaurant "Chez Matonrière".

    Dans la file, il y avait son ami monsieur Arrêter-Quand-Je-Veux.

    - Bon matin! S'écria monsieur Arrêter-Quand-Je-Veux.
}
\txtpage{<Matonrière derrière comptoir avec des cannes de haricots, avec monsieur arrêter-quand-je-veux>}

\txtpage{
    - Donc tu organises une partie de Brun? Demanda monsieur Gaboué.
    J'ai croisé madame Fin-De-Phrase en chemin et elle m'a tout expliqué.

    - Yup. Tantôt c'est un bon moment? Repris monsieur Arrêter-Quand-Je-Veux

    - Je peux pas! Y'a un trou dans mon mur.

    - Ah oui c'est vrai... Tu sais, je connais quelqu'un qui peut te faire ça pour \textit{cheap}.

    Monsieur Gaboué hocha de la tête, puis monsieur Arrêter-Quand-Je-Veux lui remis l'adresse de monsieur
    Ça-Fait-La-Job.
}
\txtpage{<Monsieur Arrêter-quand-je-veux qui remet un papier à m gaboué>}

\txtpage{
    Après un bref délai d'attente, monsieur Gaboué ressorti avec un bon bol chaud de haricots.

    Matonrière faisait vraiment les meilleurs haricots pour emporter.

    Avec tout ça il était déjà midi, alors monsieur Gaboué opta pour marcher en mangeant, question
    d'arriver plus vite chez monsieur Ça-Fait-La-Job.
}
\txtpage{<m gaboué avec une canne d'haricots pis une cuillère>}

\txtpage{
    Au coin de la rue, monsieur Gaboué tomba pile devant madame Nata.

    Monsieur Gaboué regarda madame Nata avec intention.

    Madame Nata ne bougea pas.

    Monsieur Gaboué regarda encore madame Nata, avec beaucoup d'intention cette fois-ci.

    Madame Nata ne bougea même pas d'un poil.

    Monsieur Gaboué lui fit un simple geste, puis il la contourna.
}
\txtpage{<m gaboué qui envoie un finger à mme nata>}

\txtpage{
    Monsieur Gaboué arriva enfin chez monsieur Ça-Fait-La-Job.

    - Y'a un trou dans mon mur! Annonça monsieur Gaboué.

    - Mmh. Ça arrive. Rassura monsieur Ça-Fait-La-Job. Il continua: Vous savez, je pense que j'ai
    ce qu'il vous faut.

    Puis, il alla fouiller dans son garage.
}\txtpage{<m ça-fait-la-job>}

\txtpage{
    On entendit de grands bruits. Plusieurs pots tombèrent et se cassèrent. Un chat partit en courant.

    Puis finalement, monsieur Ça-Fait-La-Job revint présenter à monsieur Gaboué une grande plaque
    de tôle.

    - Non mais je peux pas trimballer ça. Tu peux pas venir l'installer pour moi? Demanda monsieur Gaboué.

    - Oh. Je suppose que je peux appeler mes \textit{gars} pis on va te faire ça.
}
\txtpage{<tôle>}

% monsieur gaboué continue son chemin jusque chez lui, il arrive pis m ça-fait-la-job est déjà là (y'ont pris le bus)
% c'est le chaos
% dans le chaos, un mur de tôle se fait installer
% gaboué pogne les nerfs pis kickout tout le monde
% pis maintenant que tout est chill, il appelle son ami m arrêter-quand-je-veux... pour faire une partie de brun
% la musique micro-tonale c'est micro-louche

\txtpage{
    Monsieur Gaboué quitta et retourna, toujours à pied, à sa maison. Après un avant-midi pour aller chercher
    le dîner, il mit son après-midi pour faire le chemin inverse.

    Avant même qu'il puisse ouvrir sa porte, il entendit du bruit qui venait de l'intérieur. Ou plutôt \textit{du vacarme}.
}
\txtpage{<porte de m gaboué avec de la musique qui sort>}

\txtpage{
    - Oh bonjour! S'exclama monsieur Ça-Fait-La-Job lorsque la porte s'ouvrit. On a prit l'autobus.
    Comme vous n'étiez pas là on est rentré par le trou. J'espère que c'est bon?

    En plus de monsieur Ça-Fait-La-Job, il y avait monsieur Je-Suis-Pas-Fatigué, madame Je-Veux-Pas-Dormir,
    monsieur Baba, et madame Gobelin.

    Monsieur Ça-Fat-La-Job était en train de visser la tôle tout seul, tandis que les autres faisaient...
    quelque chose.
}
\txtpage{<le party chez m gaboué>}

\txtpage{
    Monsieur Gaboué, irrité, regarda monsieur Ça-Fait-La-Job avec intention.

    - Pardon. Vous verrez, j'ai presque terminé, on va pouvoir faire la fête! Répondit monsieur
    Ça-Fait-La-Job.

    Pas très rassuré, monsieur Gaboué essaya de se dire que ça serait OK.

    Madame Gobelin, qui était grimpée en haut d'une bibliothèque, remarqua la guitare de monsieur
    Gaboué.

    - Oh! On devrait faire de la musique, demanda madame Gobelin.
}
\txtpage{<madame gobelin dans une bibliothèque qui pointe une guitare>}

\txtpage{
    Madame Je-Veux-Pas-Dormir, Monsieur Je-Suis-Pas-Fatigué, et Monsieur Baba étaient très intéressés à l'idée.

    Madame Je-Veux-Pas-Dormir, alla chercher un mélodica dans son sac.

    Un pipo jaillit de nul part dans les mains de monsieur Je-Suis-Pas-Fatigué.

    Madame Gobelin saisit une poubelle qu'elle renversa, et elle s'arma de la seule cuillère de bois de
    monsieur Gaboué, ainsi que de sa seule cuillère de métal.

    Monsieur Baba prit la guitare.
}
\txtpage{<l'armement>}

\txtpage{
    Monsieur Gaboué n'aimait pas vraiment vers où les choses allaient.

    - C'est à ce moment, que monsieur Ça-Fait-La-Job s'exclamât: J'ai terminé!

    Puis il donna un coup sur le nouveau mur de tôle et fit retentir un grand GONG!

    C'était le signal: La musique se déchaîna.
}
\txtpage{<gong!>}

\txtpage{
    Il faut comprendre que la guitare de monsieur Gaboué n'avait pas été accordée, et que le mélodica
    de madame Je-Veux-Pas-Dormir, était aussi, plutôt mal accordé. Il s'en produit un effet parfaitement
    unique.

    - Le visage de madame Gobelin s'illuminât: Ça c'est de la musique microtonale!

    - Monsieur Gaboué rabroua: La musique microtonale, c'est microlouche.

    Et il en avait vraiment assez...
}
\txtpage{<capharnaum avec m gaboué qui se bouche les oreilles>}

\txtpage{
    Monsieur Gaboué leva le ton pour avoir tous leur attention, et d'un simple geste, les fit tous
    partir de chez lui.
}

\legalpage{17 novembre 2026}
\backpage

\end{document}
